%%%PAGE 188 rot=0 legibility=good summary=“原子物理和天体物理”易错点清单1-11：黑体与光子动量、能级跃迁、比结合能、开普勒定律、同点同a、衰变与半衰期、圆与椭圆轨道比较、同步卫星、氢原子光子数、行星追及求未知行星B（末尾被页底截断）
%%%BLOCK id=188-01 ch=17 sec=黑体辐射与光子动量 kind=note alt=none cont=none y=0.03-0.15
\textbf{原子物理 和 天体〉物理}% CHECK: 标题中的“〉”疑为补写“物理”的插入符

1、黑体 $T\uparrow$ $\lambda$ 极大值左移% CHECK: 原稿此处为“〉”形符号，疑为 λ（与右侧 p=h/λ 中的 λ 写法相同）

%FIG p188_a 0.34 0.062 0.47 0.135
\notefig[0.25]{figures/p188_a.png}

$p=\dfrac h\lambda$ 散射后 $\lambda\uparrow$ $p\downarrow$．
%%%END
%%%BLOCK id=188-02 ch=17 sec=氢原子能级跃迁 kind=note alt=none cont=none y=0.13-0.17
2、电子跃迁到低能级 $E_k\uparrow$ $E_{\text{电}}\downarrow$ $E_{\text{原子}}\downarrow$ 释放光子．% 注：“E原子↓”写在“E电↓”右下方、“释放光子”下方，带向下箭头
%%%END
%%%BLOCK id=188-03 ch=05 sec=卫星运行规律 kind=note alt=none cont=none y=0.13-0.16
高轨低速大周期大机大势小引力% 原稿与上一条“释放光子．”同一行的右半部分
%%%END
%%%BLOCK id=188-04 ch=17 sec=结合能与比结合能 kind=note alt=none cont=none y=0.16-0.19
3、稳定性 = 比结合能（与结合能无关）
%%%END
%%%BLOCK id=188-05 ch=05 sec=开普勒定律与中心天体 kind=note alt=none cont=none y=0.19-0.28
4、$\dfrac{r^3}{T^2}=\dfrac{GM}{4\pi^2}$ 与中心天体 $M$ 有关（绕的中心不同就不成立）

%FIG p188_b 0.17 0.222 0.275 0.27
\notefig[0.25]{figures/p188_b.png}

地球\quad $v_P=\sqrt{g_{\mathrm{I}}r_{\mathrm{I}}}\neq\sqrt{g\,r_{\text{地球}}}$% 图中小圆旁标“地球”，圆下标 Ⅰ、P
%%%END
%%%BLOCK id=188-06 ch=05 sec=卫星运行规律 kind=note alt=none cont=none y=0.28-0.32
5、同点同 $a$ 大圆大 $v$\quad 力用 $\dfrac{GMm}{r^2}$ 比．$a=\dfrac{GM}{r^2}$
%%%END
%%%BLOCK id=188-07 ch=05 sec=航天器·机械能守恒 kind=note alt=06 cont=none y=0.32-0.35
6、无动力飞行机械能守恒% 原稿此行有下划线
%%%END
%%%BLOCK id=188-08 ch=17 sec=结合能与核反应·衰变类型 kind=note alt=none cont=none y=0.34-0.41
7、生成物比结合能 大于 反应物（不是化学反应）．对放能核反应 生成物结合能也更大

生成 He 为 $\alpha$ 衰变，生成 $\mathrm{e}^-$ 为 $\beta$ 衰变，生成 $\gamma$ 为 $\gamma$ 衰变\quad 中子 $\nuc{1}{0}{n}$\quad $\nuc{0}{-1}{e}$
%%%END
%%%BLOCK id=188-09 ch=17 sec=半衰期与衰变公式 kind=note alt=none cont=none y=0.40-0.47
8、核聚变不可控．\quad 衰变公式 $\dfrac{m_{\text{余}}}{m}=\left(\dfrac12\right)^{\frac nT}$% CHECK: 左侧分数字迹潦草，疑为 m余/m；指数疑为 n/T
\quad（$n$ 为时\unclear{间}，$T$ 为半衰期）．经 $2T$ 4 个 \unclear{U}238 变为 2 个% 原稿此句上有红色斜线划过

\red{半衰期是本征属性，只与原子核有关，与外界无关}\quad \red{对大量原子核的统计规律}

（质量 \unclear{1个} 减 1 半 $\times$，不能对少量物质说，只能广泛地说）
%%%END
%%%BLOCK id=188-10 ch=05 sec=卫星变轨·圆轨道与椭圆轨道比较 kind=note alt=none cont=none y=0.46-0.61
%FIG p188_c 0.05 0.458 0.27 0.568
\notefig[0.25]{figures/p188_c.png}

$r_{\text{甲}}=r_{\text{乙}}$% CHECK: “乙”手写近似“2”，图中椭圆右端亦有类似“乙”的标记

近地点 $a_{\text{乙}}>a_{\text{甲}}$\quad\quad $a=\dfrac{\Delta v}{\Delta t}$ or $\dfrac Fm$\quad $a$ 与卫星质量无关．

由开三 $T_{\text{甲}}=T_{\text{乙}}$

相同时间内扫过面积不同

构造的轨道．（其上方有一箭头指向图中椭圆下方）
%%%END
%%%BLOCK id=188-11 ch=17 sec=核反应方程·β衰变 kind=formula alt=none cont=none y=0.63-0.66
$\beta$ 衰变\quad $\nuc{1}{0}{n}\to\nuc{1}{1}{H}+\nuc{0}{-1}{e}$
%%%END
%%%BLOCK id=188-12 ch=05 sec=同步卫星 kind=note alt=none cont=none y=0.66-0.76
9、$V_{\text{倾斜轨}}=V_{\text{同步}}$ 但动能 $\dfrac12mv^2$ 要比动能要有质量关系．% 原稿“要比动能要有质量关系”下有下划线

$V_{\text{同步}}>V_{\text{月绕地}}$\quad 同步轨道 $3.6\times10^4\unit{km}$\quad 地球半径 $6.4\times10^3\unit{km}$

$T_{\text{同步}}=30\unit{h}$\quad 高轨低速大周期大机大势小引力% CHECK: 同步卫星周期应为 24h，原稿疑写作 30（或 24 误辨）
%%%END
%%%BLOCK id=188-13 ch=17 sec=氢原子能级·光子与光电效应 kind=note alt=none cont=none y=0.76-0.84
10、一群氢原子有 $\dfrac{n(n-1)}{2}$ 种光子\quad 1 个氢原子有 $n-1$ 种光子．（光电效应仅表粒子性）．

$E_k=h\nu-W_0$\quad $p=\dfrac h\lambda$\quad $p=\dfrac Ec$\quad 故光子能量越大 动量也越大．\quad $\lambda=\dfrac hp=\dfrac{h}{\sqrt{2mE}}$．
%%%END
%%%BLOCK id=188-14 ch=05 sec=天体追及·未知行星B kind=problem alt=none cont=none y=0.84-1.00
\begin{problem}[11]
A 每 $t_0$ 发生一次最大偏离．（A 外有一颗未知行星 B），与 A 同方向绕行 求 B 参数．

%FIG p188_d 0.02 0.83 0.21 0.925
\notefig[0.25]{figures/p188_d.png}

\begin{solution}
$\dfrac{t_0}{T_A}-\dfrac{t_0}{T_B}=1$．\quad $T_B=\dfrac{T_At_0}{t_0-T_A}$\quad $\dfrac{GMm'}{R_B^{2}}=m'\left(\dfrac{2\pi}{T_B}\right)^{2}R_B$．

$R_B=\sqrt[3]{\left(\dfrac{t_0}{t_0-T_A}\right)^{2}}$% CHECK: 根号前疑漏 R_0

$V_B=\dfrac{2\pi R_B}{T_B}=\dfrac{2\pi R_0}{T_0}\sqrt[3]{\dfrac{t_0-T_A}{\text{\unclear{…}}}}$\quad $a_B=\dfrac{V_B^{2}}{R_B}=\dfrac{4\pi^{2}R_0}{T_0^{2}}\sqrt[3]{\dfrac{(t_0-T_A)^{4}}{t_0}}$% CHECK: 页底被截断，两处根号内分母不完整（第一处只见残笔，第二处只见 t_0，指数可能被截）
\end{solution}
\end{problem}
%%%END
%%%PAGEEND 188
